\documentclass[11pt]{article}

\usepackage[T1]{fontenc}
\usepackage{lmodern}
\usepackage[margin=1in]{geometry}
\usepackage{amsmath,amssymb,amsthm,mathtools}
\usepackage{microtype}
\usepackage{enumitem}
\usepackage{xcolor}
\usepackage{graphicx}
\usepackage{hyperref}
\usepackage[nameinlink,noabbrev]{cleveref}

\hypersetup{
  colorlinks=true,
  linkcolor=blue!55!black,
  citecolor=blue!55!black,
  urlcolor=blue!55!black,
  pdftitle={An Exact Obstruction to Joint Sequential Continuity of Abstract Measure Addition},
  pdfauthor={Run fa_banach_001, lane 2}
}

\newtheorem{theorem}{Theorem}[section]
\newtheorem{lemma}[theorem]{Lemma}
\newtheorem{corollary}[theorem]{Corollary}
\newtheorem{proposition}[theorem]{Proposition}
\theoremstyle{definition}
\newtheorem{definition}[theorem]{Definition}
\theoremstyle{remark}
\newtheorem{remark}[theorem]{Remark}

\newcommand{\N}{\mathbb N}
\newcommand{\C}{\mathbb C}
\newcommand{\ellone}{\ell_1}
\newcommand{\taumu}{\tau_\mu}
\newcommand{\wstar}{w^*}

\title{\textbf{An Exact Obstruction to Joint Sequential Continuity}\\
\large Abstract measure topology on L-embedded Banach spaces}
\author{Run \texttt{fa\_banach\_001}, lane 2}
\date{19 July 2026}

\begin{document}
\maketitle

\begin{abstract}
Peralta and Pfitzner ask whether addition is sequentially jointly continuous
for the abstract measure topology on every L-embedded Banach space.  We give
an exact obstruction reduction.  If addition fails, two seminormalized
asymptotically isometric copies of the unit vector basis of $\ellone$ can be
chosen whose termwise sum is seminormalized and either (i) is an $\ellone$-
basic sequence with no almost isometric $\ellone$ subsequence, exhibiting
strict norm loss at every singular weak-star cluster point, or (ii) is weakly
null.  Conversely, either configuration forces failure.  In the weakly null
branch the two summands have the same limiting norm and opposite singular
weak-star limits.  A simple non-L-embedded max-sum realizes the second branch,
showing that the remaining issue is genuinely tied to L-embedded geometry.
The universal question is not resolved here.
\end{abstract}

\section*{Packet status}

\textbf{partial\_result\_likely\_valid.}  The reduction and its converse are
proved.  This packet does not claim that the two obstructions are impossible,
or that either occurs, in an L-embedded Banach space.

\tableofcontents

\section{Source question}

Let $X$ be L-embedded, with its canonical L-decomposition
\[
             X^{**}=X\oplus_1 X_s.
\]
Pfitzner defined the abstract measure topology $\taumu$ by prescribing its
convergent sequences \cite{Pfitzner}.  Peralta and Pfitzner recall the
definition and then ask the following on page 6 of arXiv:1505.05302
\cite{PeraltaPfitzner}.

\begin{quote}
By definition addition is separately continuous with respect to $\taumu$,
but it is an open problem whether addition is sequentially jointly continuous
or even jointly continuous on an L-embedded space $X$.
\end{quote}

\begin{figure}[htbp]
\centering
\includegraphics[width=\textwidth]{figures/open_problem_crop.png}
\caption{The source question and the sequential description of $\taumu$,
page 6 of arXiv:1505.05302.}
\end{figure}

The same paper proves sequential joint continuity when $X$ is a
JBW$^*$-triple predual.  The question for arbitrary L-embedded spaces remains
the target of this packet.

\section{Definitions and preliminary facts}

\begin{definition}
A sequence $(u_n)$ of nonzero vectors \emph{spans $\ellone$
asymptotically isometrically} if there are $\delta_n\in[0,1)$ with
$\delta_n\to0$ such that, for every finitely supported scalar sequence
$(a_n)$,
\[
 \sum_n(1-\delta_n)|a_n|
 \leq \left\|\sum_n a_n\frac{u_n}{\|u_n\|}\right\|
 \leq \sum_n|a_n|.
\]
It is \emph{seminormalized} if
$0<\inf_n\|u_n\|\leq\sup_n\|u_n\|<\infty$.
\end{definition}

For an L-embedded space, a sequence is $\taumu$-null exactly when it is
bounded and every subsequence has a further subsequence which is norm null or
whose normalization spans $\ellone$ almost isometrically
\cite{Pfitzner,PeraltaPfitzner}.  Pfitzner's perturbation result allows
``almost isometrically'' to be replaced, after one more subsequence, by
``asymptotically isometrically.''  We use that equivalent sequential
description below.

We also use three standard facts.
\begin{enumerate}[label=(\roman*)]
\item Every weak-star accumulation point of an asymptotically isometric
$\ellone$ sequence in an L-embedded space lies in $X_s$, by Pfitzner's
fundamental weak-star accumulation lemma \cite{Pfitzner}.
\item A bounded sequence has either an $\ellone$-basic subsequence or a weakly
Cauchy subsequence, by Rosenthal's $\ellone$ theorem \cite{Rosenthal}.
\item L-embedded Banach spaces are weakly sequentially complete
\cite[Theorem IV.2.2]{HWW}.
\end{enumerate}

The following norm refinement of (i) is elementary but useful.

\begin{lemma}[No norm loss for an asymptotic summand]
\label{lem:no-loss}
Let $(u_n)$ span $\ellone$ asymptotically isometrically in an L-embedded
space, and suppose $\|u_n\|\to a>0$.  Every weak-star accumulation point
$u^{**}$ of $(u_n)$ belongs to $X_s$ and satisfies $\|u^{**}\|=a$.
\end{lemma}

\begin{proof}
Membership in $X_s$ is Pfitzner's result recalled above.  Put
$q_n=u_n/\|u_n\|$.  If a subnet of $(u_n)$ converges weak-star to $u^{**}$,
the corresponding subnet of $(q_n)$ converges to $q^{**}=u^{**}/a$.

Fix $\varepsilon>0$ and choose $N$ so that $\delta_n\leq\varepsilon$ for
$n\geq N$.  On the closed span of $(q_n)_{n\geq N}$, the formula
\[
       f_N\left(\sum_{n\geq N}c_nq_n\right)=\sum_{n\geq N}c_n
\]
defines a functional of norm at most $(1-\varepsilon)^{-1}$.  Extend it to
$X$ by Hahn--Banach.  Along any free subnet of the tail,
$f_N(q_n)=1$, and hence $f_N(q^{**})=1$.  Thus
$\|q^{**}\|\geq1-\varepsilon$.  Since $\|q^{**}\|\leq1$ and
$\varepsilon$ is arbitrary, $\|q^{**}\|=1$.
\end{proof}

\section{Proof intuition}

Failure of joint sequential continuity means that two $\taumu$-null
sequences have a sum which fails the subsequence test defining $\taumu$.
After discarding the easy norm-null cases, both summands must be
seminormalized asymptotically isometric $\ellone$ sequences.  Their weak-star
cluster points therefore live, with full norm, in the singular summand $X_s$.
Consequently every weak-star cluster point of their termwise sum also lies in
$X_s$.

Rosenthal's dichotomy now leaves only two possibilities for a bad sum.  It
either retains an $\ellone$-basic subsequence but never improves to an almost
isometric one; this is the singular norm-loss branch.  Or it has no
$\ellone$ subsequence, hence has a weakly convergent subsequence.  Its weak
limit belongs simultaneously to $X$ and $X_s$, so it must be zero.  This is
the weak-null cancellation branch.  Reversing the argument is immediate:
neither branch can pass the defining $\taumu$ subsequence test.

\section{The exact obstruction theorem}

\begin{theorem}[Two-branch obstruction]
\label{thm:obstruction}
Let $X$ be an L-embedded Banach space with abstract measure topology
$\taumu$.  Addition fails to be sequentially jointly continuous if and only
if there are seminormalized sequences $(x_n)$ and $(y_n)$, each spanning
$\ellone$ asymptotically isometrically, such that
$z_n=x_n+y_n$ is seminormalized and, after passing to a subsequence, one of
the following mutually exclusive alternatives holds.

\begin{enumerate}[label=\textup{(\Alph*)}]
\item \textbf{Distorted-$\ellone$ / singular norm-loss branch.}
$(z_n)$ is equivalent to the unit vector basis of $\ellone$, but no
subsequence of $(z_n/\|z_n\|)$ spans $\ellone$ almost isometrically.
Equivalently via Pfitzner's cluster criterion, every weak-star accumulation
point of $(z_n/\|z_n\|)$ lies in $X_s$ and has norm strictly less than one.

\item \textbf{Weak-null cancellation branch.}
$(z_n)$ is weakly null.  Moreover, the sequences may be chosen so that
$\|x_n\|\to a$, $\|y_n\|\to a$ for some $a>0$, and every common weak-star
cluster pair $(x_s,y_s)$ satisfies $x_s,y_s\in X_s$ and $y_s=-x_s$.
\end{enumerate}
\end{theorem}

\begin{proof}
Suppose first that addition is not sequentially jointly continuous.  There
are $\taumu$-null sequences $(a_n)$ and $(b_n)$ such that
$c_n=a_n+b_n$ is not $\taumu$-null.  By the defining subsequence criterion,
there is a subsequence of $(c_n)$ with no further subsequence which is norm
null or whose normalization spans $\ellone$ almost isometrically.  Relabel
this subsequence as $(c_n)$.

Apply the $\taumu$-null criterion first to $(a_n)$ and then, inside the
resulting subsequence, to $(b_n)$.  A norm-null alternative for either
summand is impossible: if the other summand is norm null then so is $(c_n)$;
if the other is an almost/asymptotically isometric $\ellone$ sequence, adding
a norm-null sequence preserves that property after separating the norms away
from zero.  Either conclusion contradicts the choice of $(c_n)$.  Passing to
further subsequences and using Pfitzner's almost-to-asymptotic extraction, we
therefore obtain seminormalized asymptotically isometric $\ellone$ sequences
$(x_n)$ and $(y_n)$.  We may also arrange
\[
       \|x_n\|\to a>0,\qquad \|y_n\|\to b>0,
\]
and $z_n=x_n+y_n$ remains seminormalized and has no good subsequence.

Every weak-star accumulation point of $(z_n)$ lies in $X_s$.  Indeed, given
a weak-star convergent subnet of $(z_n)$, pass to a further subnet on which
the bounded pair $(x_n,y_n)$ converges weak-star to $(x_s,y_s)$.  By
\cref{lem:no-loss}, $x_s,y_s\in X_s$, and hence the limit
$x_s+y_s$ of the sum belongs to $X_s$.

If $(z_n)$ has an $\ellone$-basic subsequence, pass to it.  It still has no
almost isometric $\ellone$ subsequence.  Pfitzner's cluster criterion says
that a norm-one weak-star accumulation point in $X_s$ would produce an
asymptotically, hence almost, isometric $\ellone$ subsequence.  Thus every
normalized cluster point has norm less than one, and alternative (A) holds.

Suppose instead that $(z_n)$ has no $\ellone$-basic subsequence.  Rosenthal's
theorem gives a weakly Cauchy subsequence.  Since $X$ is weakly sequentially
complete, pass to a weakly convergent subsequence, say $z_n\rightharpoonup z$
with $z\in X$.  The same $z$, viewed in $X^{**}$, is a weak-star accumulation
point of $(z_n)$, so the preceding paragraph gives $z\in X_s$.  The
L-decomposition is direct, hence $z\in X\cap X_s=\{0\}$.  Thus $(z_n)$ is
weakly null and alternative (B) holds.

For the final assertion in (B), take a common weak-star cluster pair
$(x_s,y_s)$.  Weak convergence of the sum gives $x_s+y_s=0$.  By
\cref{lem:no-loss}, $\|x_s\|=a$ and $\|y_s\|=b$, so $a=b$ and
$y_s=-x_s$.

Conversely, assume either (A) or (B).  Each of $(x_n)$ and $(y_n)$ is
$\taumu$-null.  In (A), $(z_n)$ is seminormalized and has no almost
isometric $\ellone$ subsequence, so it is not $\taumu$-null.  In (B), a
weakly null sequence cannot have an $\ellone$-basic subsequence: on the span
of such a subsequence the functional taking value one on every basis vector
is bounded and extends to $X$, contradicting weak nullity.  Since $(z_n)$ is
also seminormalized, it has neither kind of good subsequence and is not
$\taumu$-null.  Addition therefore fails to be sequentially jointly
continuous.
\end{proof}

\begin{corollary}[Operator form of branch (B)]
Branch \textup{(B)} occurs if and only if there are asymptotically isometric
embeddings $U,V:\ellone\to X$ such that
\[
       (U+V)e_n \rightharpoonup 0,
       \qquad \inf_n\|(U+V)e_n\|>0,
\]
after harmless positive rescaling of the two bases by factors converging to
the same constant.
\end{corollary}

\begin{proof}
In branch (B), write $\|x_n\|\to a$ and $\|y_n\|\to a$, and set
$q_n=x_n/\|x_n\|$, $r_n=y_n/\|y_n\|$.  Then
\[
 q_n+r_n=\frac{z_n}{a}
 +\left(\frac{1}{\|x_n\|}-\frac1a\right)x_n
 +\left(\frac{1}{\|y_n\|}-\frac1a\right)y_n.
\]
The last two terms tend to zero in norm.  Hence $(q_n+r_n)$ is weakly null
and, after deleting finitely many terms, seminormalized.  Define
$Ue_n=q_n$ and $Ve_n=r_n$.  The asymptotic lower estimates make $U$ and
$V$ embeddings.  The converse follows directly from the unit-vector images.
\end{proof}

This operator formulation may be useful because the operator $U+V$ is weakly
compact: the weakly null set
$\{(U+V)e_n:n\in\N\}\cup\{0\}$ is weakly compact, and its closed absolutely
convex hull contains the image of the unit ball of $\ellone$.

\section{A sharpness model outside the L-embedded class}

The weak-null obstruction is geometrically feasible without L-embeddedness.
Let
\[
            E=\ellone\oplus_\infty c_0
\]
and write $(e_n)$ for the unit vector basis in either coordinate.  Define
\[
       x_n=(e_n,0),\qquad y_n=(-e_n,e_n),\qquad z_n=x_n+y_n=(0,e_n).
\]
Both $(x_n)$ and $(y_n)$ span $\ellone$ isometrically.  Indeed, for finitely
supported $(a_n)$,
\[
 \left\|\sum_n a_ny_n\right\|
 =\max\left\{\sum_n|a_n|,\sup_n|a_n|\right\}
 =\sum_n|a_n|.
\]
Yet $(z_n)$ is weakly null and has norm one.  Thus it is exactly the branch
(B) configuration.

This is not a counterexample to the source question.  The space $E$ is not
L-embedded: L-embedded spaces are weakly sequentially complete, while the
$c_0$ summand contains the weakly Cauchy sequence
$e_1+\cdots+e_n$, which has no weak limit in $c_0$.

\section{Verification, limitations, and novelty}

No computation is used.  The proof should be reviewed at three points:
\begin{enumerate}[label=(\arabic*)]
\item the double subsequence extraction from the definition of $\taumu$;
\item the passage from individual singular cluster points to singular cluster
points of the termwise sum;
\item the use of Rosenthal's theorem and weak sequential completeness to turn
the non-$\ellone$ branch into a weakly null sequence.
\end{enumerate}

The theorem does not exclude either obstruction in an arbitrary L-embedded
space.  In particular, it is not a full solution of the source question.
Its contribution is an exact reduction to two concrete sequence-geometric
configurations, including a complete description of the cancellation in the
weak-null branch.

A bounded literature search through 19 July 2026 used the exact phrases
``abstract measure topology'' with ``addition'', ``sequentially jointly
continuous'', and ``weakly null'', as well as searches for termwise sums of
asymptotically isometric $\ellone$ sequences.  It found the source papers and
later applications of abstract measure topology, but no later proof,
counterexample, or this two-branch reduction.  This is evidence for review,
not a claim of priority or an exhaustive citation search.

\section{Human-review recommendation}

Review as a likely valid substantial partial reduction.  The most promising
next target is branch (B): determine whether an L-embedded space can contain
two asymptotically isometric $\ellone$ embeddings whose sum is weakly compact
but remains seminormalized on the unit-vector basis.  Branch (A) is the
secondary target and asks whether the singular cluster points of such a sum
can all lose norm.

\begin{thebibliography}{9}

\bibitem{PeraltaPfitzner}
A.~M. Peralta and H.~Pfitzner,
\emph{Weak Banach--Saks property and Koml\'os' theorem for preduals of
JBW$^*$-triples}, arXiv:1505.05302, 2015.

\bibitem{Pfitzner}
H.~Pfitzner,
\emph{L-embedded Banach spaces and measure topology},
Israel J. Math. \textbf{205} (2015), 421--451;
arXiv:math/0003154.

\bibitem{Rosenthal}
H.~P. Rosenthal,
\emph{A characterization of Banach spaces containing $\ell^1$},
Proc. Nat. Acad. Sci. USA \textbf{71} (1974), 2411--2413.

\bibitem{HWW}
P.~Harmand, D.~Werner, and W.~Werner,
\emph{M-Ideals in Banach Spaces and Banach Algebras},
Lecture Notes in Mathematics 1547, Springer, 1993.

\end{thebibliography}

\end{document}
